\centerline{\bf \Huge Вступительный разнобой}

\begin{flushright}
        \scriptsize{хочу спать}
\end{flushright}

\problem{1} Пусть $p_1, p_2, \ldots, p_{100}$ - сто простых чисел, среди которых нет одинаковых. Натуральные числа $a_1, \ldots, a_k$, большие 1 , таковы, что каждое из чисел $p_1 p_2^3, p_2 p_3^3, \ldots, p_{99} p_{100}^3, p_{100} p_1^3$ равно произведению каких-то двух из чисел $a_1, \ldots, a_k$. Докажите, что $k \geq 150$. 

\problem{2}
Будем называть натуральное число красивым, если в его десятичной записи поровну цифр 0,1,2, а других цифр нет. Может ли произведение двух красивых чисел быть красивым?

\problem{3}
Числа $x$ и $y$, не равные 0, удовлетворяют неравенствам $x^2-x>y^2 u y^2-y>x^2$. Какой знак может иметь произведение $xy$ (укажите все возможности)?

\problem{4}
Дано натуральное число n, большее 2. Докажите, что если число $n!+n^3+1-$ простое, то число $n^2+2$ представляется в виде суммы двух простых чисел.

\problem{5}
Докажите, что существует натуральное число $n$, большее $10^{100}$, такое, что сумма всех простых чисел, меньших n, взаимно проста с n.
